\chapter{Environmental Assessment}\label{sec:EnvironmentalAssessment}

This chapter assesses the environmental impact of carrying out this Master's 
thesis, expressed in terms of greenhouse gas emissions, \ac{CO2e}, resulting 
from the electricity consumed during the experimental and computational work.

\section{Electricity Consumption}

Three sources of electricity consumption were considered: the energy drawn from 
the grid to recharge the cell, the operating power of the laboratory instruments 
during the tests, and the power consumed by the student's laptop 
throughout the project.

\paragraph{Cell recharging.} The cell was fully recharged approximately 15 times 
over the course of the characterisation and validation tests, drawing 
approximately $111.4\,\text{Wh}$ per charge, as measured directly during the 
charge test (Equation~\ref{eq:EChargeDiscrete}):

\begin{equation}
    15 \times 111.4\,\text{Wh} = 1.671\,\text{kWh}
\end{equation}

\paragraph{Laboratory instruments.} The power analyser, cycler, and current clamp 
were in active use for a total of approximately $70.5$ hours, distributed across 
the four test as summarised in Table~\ref{tab:InstrumentHours}, at an 
assumed operating power of $20\,\text{W}$:

\begin{equation}
    70.5\,\text{h} \times 20\,\text{W} = 1.410\,\text{kWh}
\end{equation}

\begin{table}[H]
    \centering
    \begin{tabular}{lr}
        \toprule
        \textbf{Test campaign} & \textbf{Hours} \\
        \midrule
        Charge tests (15 sessions $\times$ 2~h)                  & 30.0 \\
        Discharge tests (2 sessions $\times$ 7~h)                & 14.0 \\
        Pulse discharge tests (2 sessions $\times$ 5~h)          & 10.0 \\
        Random step-current tests (2 full $\times$ 6~h, 3 trial $\times$ 1.5~h) & 16.5 \\
        \midrule
        \textbf{Total} & \textbf{70.5} \\
        \bottomrule
    \end{tabular}
    \caption{Estimated instrument usage hours by test campaign.}
    \label{tab:InstrumentHours}
\end{table}

\paragraph{Personal computer.} The student's laptop was used for approximately 
$500$ hours over the course of the project on MATLAB simulations for data 
processing, and thesis writing at an average power draw of $30\,\text{W}$. 
Approximately two weeks were spent working from home during the 
literature review phase, with the remaining time split approximately $30\%$ at public 
libraries and $70\%$ in the laboratory.

\begin{equation}
    500\,\text{h} \times 30\,\text{W} = 15.000\,\text{kWh}
\end{equation}

\begin{table}[H]
    \centering
    \begin{tabular}{lr}
        \toprule
        \textbf{Source} & \textbf{Energy [kWh]} \\
        \midrule
        Cell recharging                & 1.671 \\
        Laboratory instruments         & 1.410 \\
        Personal laptop                & 15.000 \\
        \midrule
        \textbf{Total}                 & \textbf{18.081} \\
        \bottomrule
    \end{tabular}
    \caption{Total estimated electricity consumption.}
    \label{tab:TotalElectricityEnv}
\end{table}

\section{Greenhouse Gas Emissions}

The Spanish national electricity grid mix emission factor for 2025 
is $0.258\,\text{kg}\,\ac{CO2e}/\text{kWh}$, 
consistent with the emission factors presented by the Catalan Office for Climate 
Change~\cite{canviclimatic}. Applying this factor to the total electricity 
consumption:

\begin{equation}
    18.081\,\text{kWh} \times 0.258\,\text{kg}\,\ac{CO2e}/\text{kWh} = 4.665\,\text{kg}\,\ac{CO2e}
\end{equation}

\section{Environmental Impact Discussion}

The estimated carbon footprint of this thesis, $4.665\,\text{kg}\,\ac{CO2e}$. 
The personal laptop accounts for approximately $83\%$ of the total electricity 
consumption. The cell 
recharging and instrument operation together account for less than $20\%$ of the 
total footprint, reflecting the efficiency of the algebraic parameter 
identification method used in this work (Section~\ref{Parameters}), which 
requires substantially shorter tests, and therefore less instrument energy 
consumption than conventional characterisation approaches requiring long 
relaxation periods.


\chapter{Social and Gender Equality Assessment}\label{sec:SocialGenderAssessment}

\section{Access to the Technology}

The battery state estimation framework developed in this thesis is not a product or service with differentiated 
access. It does not discriminate between people of different gender identities or against vulnerable 
groups. The algorithms, once published, are equally accessible to any engineer or 
researcher regardless of gender. 

The access to electric vehicle technology, of which battery management is a component, remains unevenly distributed 
across socio-economic groups due to the current high cost of electric vehicles. This 
is a systemic issue of the electric mobility sector.

\section{Differentiated Impacts of the Research}

The results of this thesis, \ac{SOC} and \ac{SOH} estimation 
accuracy for a racing motorcycle battery, do not have any known differentiated 
impact on any gender, either directly or indirectly.

\section{Gender Distribution in Cited Sources}
This section examines the gender distribution 
among the authors of the works cited throughout this document, as a reflection of 
representation within the research field.

This analysis was performed by inferring the likely gender associated with each 
author's given name. This method has clear limitations: it can misclassify 
individuals whose name does not conventionally align with their gender identity 
and many given names, particularly transliterated names common in East Asian 
publications, do not carry a strong gender association for a non-native reader 
and were therefore classified as ambiguous rather than guessed. Organisations, 
standards bodies, and manufacturer datasheets cited without an individual named 
author were excluded from this count, as were entries not ultimately used in this 
thesis.

Of the 51 individually named authors across the references cited in this thesis, 
24 ($47.1\%$) were classified as having a male-associated given name, 4 ($7.8\%$) 
as female-associated, and 23 ($45.1\%$) as ambiguous. Restricting the analysis to 
only the 28 names with a reasonably inferable gender, $85.7\%$ are male-associated 
and $14.3\%$ are female-associated.

This reflects a 
research field in which women remain a clear minority among identifiable authors, 
independent of any specific choice made during the literature review process.

\section{Relevant Sustainable Development Goals}

This work contributes to the following \ac{SDG}:

\begin{itemize}
    \item \textbf{SDG 7 -- Affordable and Clean Energy}: accurate battery state 
    estimation improves the efficiency and safe utilisation of \ac{Li-ion} 
    battery capacity, supporting the broader adoption of electric mobility.
    \item \textbf{SDG 9 -- Industry, Innovation and Infrastructure}: this work 
    contributes methodological innovation in \ac{BMS} technology, an enabling 
    infrastructure for electric vehicles.
    \item \textbf{SDG 13 -- Climate Action}: by supporting the reliability and 
    performance of electric vehicles, this work contributes indirectly to the 
    decarbonisation of the transport sector, as discussed in the 
    Preface.
\end{itemize}